\chapter{Συμπεράσματα και μελλοντικές επεκτάσεις}
\label{ch:symperasmata}

\section{Σύνοψη}
\label{sec:synopsi}

Η εργασία ξεκίνησε από μια παρατήρηση για τις αγορές \en{NFT}: είναι οργανωμένες γύρω από
το νόμισμα, και ο χρήστης που θέλει να ανταλλάξει ένα αντικείμενο με ένα άλλο, ιδίως
διαφορετικής συλλογής, δεν εξυπηρετείται. Το \en{Sui}, του οποίου το μοντέλο αντικειμένων
προσφέρεται για ασφαλή δέσμευση χωρίς διαρκείς εγκρίσεις προς συμβόλαια, δεν διέθετε
τέτοια λειτουργία.

Σχεδιάστηκε ένα πρωτόκολλο στο οποίο η προσφορά είναι η ίδια ένα αντικείμενο της αλυσίδας
που κρατά δεσμευμένα τα προσφερόμενα στοιχεία, με τέσσερις πράξεις: υποβολή, αποδοχή,
απόρριψη και ακύρωση. Το πρωτόκολλο υλοποιήθηκε σε δύο \en{modules} \en{Move} συνολικού
μήκους 323 γραμμών, με γενικούς τύπους για τις συλλογές και δυναμικά πεδία αντικειμένων
για τη δέσμευση, γραμμένα σύμφωνα με δημοσιευμένα πρότυπα ύφους και αρχιτεκτονικής για
\en{Sui Move}, και συνοδεύεται από 21 ελέγχους μονάδας. Αναπτύχθηκε εφαρμογή ιστού που
λειτουργεί χωρίς διακομιστή, ανακαλύπτει \en{NFTs} και προσφορές από δεδομένα της
αλυσίδας, και παρουσιάζει κάθε προσφορά ως ανταλλαγή «δίνεις \swap{} παίρνεις». Η εφαρμογή
επαληθεύθηκε από άκρη σε άκρη σε τρία επίπεδα: ελέγχους μονάδας, 20 σενάρια συμβολαίου σε τοπικό
δίκτυο και 17 ροές αυτοματοποιημένης πλοήγησης.

Οι δοκιμές αντοχής έδειξαν ότι το κόστος μιας προσφοράς κλιμακώνεται γραμμικά έως
εκατοντάδες δεσμευμένα \en{NFTs}, ότι εκατό ταυτόχρονες ανταλλαγές περνούν χωρίς καμία
αποτυχία, και ότι είκοσι ταυτόχρονες αποδοχές της ίδιας προσφοράς καταλήγουν πάντοτε σε
ακριβώς μία επιτυχία. Η προβολή στο \en{mainnet}, με τις τρέχουσες τιμές του δικτύου,
φέρνει το κόστος μιας πλήρους ανταλλαγής στα 0{,}0002 \en{SUI}.

\section{Συμπεράσματα}
\label{sec:conclusions}

\paragraph{Το μοντέλο αντικειμένων απλοποιεί ουσιαστικά το πρόβλημα.} Η δέσμευση
στοιχείων, που σε μοντέλο λογαριασμών απαιτεί εγκρίσεις και προσεκτική λογιστική, στο
\en{Sui} είναι μια μεταβίβαση κυριότητας σε αντικείμενο. Η εξουσιοδότηση της απόρριψης
υλοποιήθηκε χωρίς κανένα μητρώο κατόχων, απαιτώντας απλώς μια αναφορά στο αντικείμενο. Η
ορθότητα υπό ταυτόχρονη πρόσβαση δεν απαίτησε κώδικα: προκύπτει από το ότι η προσφορά
είναι κοινόχρηστο αντικείμενο που διαγράφεται στην πρώτη επιτυχή ενέργεια, και το σενάριο
των δύο ταυτόχρονων αποδοχών το επιβεβαίωσε στην πράξη. Η ίδια ιδιότητα έχει μια συνέπεια
που αξίζει να γίνεται ρητή στους χρήστες: η προσφορά δεσμεύεται στο αντικείμενο, όχι στον
κάτοχό του, και τον ακολουθεί αν αυτό αλλάξει χέρια.

\paragraph{Το σύστημα τύπων της \en{Move} λειτουργεί ως έλεγχος ορθότητας.} Τα πιο
επικίνδυνα σφάλματα ενός πρωτοκόλλου δέσμευσης, η απώλεια ή ο εγκλωβισμός στοιχείων, είναι
σε μεγάλο βαθμό σφάλματα μεταγλώττισης. Ένα \en{NFT} ή ένα υπόλοιπο που δεν διατίθεται
ρητά εμποδίζει τη μεταγλώττιση. Το ένα σημείο όπου η γλώσσα δεν βοηθά, η διαγραφή
αντικειμένου με εναπομείναντα δυναμικά πεδία, ήταν ακριβώς εκείνο που χρειάστηκε
επιχειρηματολογία και ειδικούς ελέγχους.

\paragraph{Οι γενικοί τύποι αρκούν για τις διαφορετικές συλλογές.} Η υποστήριξη
αυθαίρετων συλλογών δεν απαίτησε μητρώο, διαχειριστή ή ενδιάμεσο \en{token}. Δύο
παράμετροι τύπου στη δομή της προσφοράς μεταφέρουν τον έλεγχο στην εικονική μηχανή. Το
τίμημα είναι ο περιορισμός σε έναν τύπο προσφερομένων ανά προσφορά.

\paragraph{Μια εφαρμογή χωρίς διακομιστή είναι εφικτή, με όρια.} Τα συμβάντα επαρκούν για
να ανακαλύψει η εφαρμογή ό,τι χρειάζεται από έναν πλήρη κόμβο. Η προσέγγιση αυτή δίνει
ισχυρές εγγυήσεις διαφάνειας και μηδενικό κόστος λειτουργίας, αλλά δεν κλιμακώνεται, διότι
κάθε ερώτημα διατρέχει όλο το ιστορικό.

\paragraph{Το κόστος είναι χαμηλό και ανεξάρτητο από την αξία.} Μια πλήρης ανταλλαγή
κοστίζει περίπου 0{,}002 \en{SUI} σε \en{gas}, χωρίς τέλη. Η μέτρηση ανέδειξε και δύο
πράγματα που δεν είχαν προβλεφθεί: η εγγύηση αποθήκευσης επιστρέφεται σε όποιον κλείνει την
προσφορά και όχι σε όποιον την κατέβαλε, και κάθε πλεονάζον πεδίο σε ένα αντικείμενο
κοστίζει. Η αφαίρεση δύο πεδίων που επαναλάμβαναν πληροφορία διαθέσιμη αλλού μείωσε το
κόστος αποθήκευσης της προσφοράς κατά σχεδόν ένα τρίτο.

\paragraph{Η επαλήθευση από άκρη σε άκρη αποκαλύπτει ό,τι ο έλεγχος τύπων δεν μπορεί.}
Το σοβαρότερο πρόβλημα που συναντήθηκε στην εφαρμογή, η απουσία μιας μεθόδου από τον
τοπικό κόμβο, δεν προκαλούσε κανένα σφάλμα μεταγλώττισης και εντοπίστηκε μόνο όταν οι ροές
εκτελέστηκαν σε πραγματικό δίκτυο. Η επένδυση σε ένα προσωρινό πορτοφόλι και σε
αυτοματοποιημένη πλοήγηση αποδόθηκε πολλαπλά: εντόπισε σφάλματα, και παρήγαγε τα
στιγμιότυπα της εργασίας.

\section{Μελλοντικές επεκτάσεις}
\label{sec:future}

Το σύστημα μπορεί να επεκταθεί προς πολλές κατευθύνσεις. Παρατίθενται κατά σειρά
εγγύτητας προς την τρέχουσα υλοποίηση.

\begin{enumerate}
  \item \textbf{Λήξη προσφορών.} Προσθήκη προαιρετικής χρονικής σήμανσης λήξης, που
        ελέγχεται κατά την αποδοχή με το αντικείμενο \code{Clock} του συστήματος. Μια
        ληγμένη προσφορά θα μπορούσε να κλείνεται από οποιονδήποτε, με επιστροφή των
        στοιχείων στον προσφέροντα.
  \item \textbf{Απόδοση της εγγύησης στον προσφέροντα.} Το συμβόλαιο θα μπορούσε να
        δέχεται κατά την υποβολή ένα μικρό ποσό που επιστρέφεται στον προσφέροντα κατά τον
        τερματισμό, αντισταθμίζοντας την επιστροφή αποθήκευσης που εισπράττει ο κάτοχος.
        Εναλλακτικά, η διεπαφή θα μπορούσε απλώς να κάνει το ποσό ορατό.
  \item \textbf{Αντιπροτάσεις.} Ο κάτοχος να μπορεί να απαντήσει σε μια προσφορά με
        τροποποιημένους όρους. Υλοποιείται ως νέα προσφορά με αντίστροφους ρόλους που
        παραπέμπει στην αρχική, ώστε η διεπαφή να παρουσιάζει τη διαπραγμάτευση ως ενιαίο
        νήμα.
  \item \textbf{Ετερογενή καλάθια.} Άρση του περιορισμού ενός τύπου προσφερομένων ανά
        προσφορά. Τα δυναμικά πεδία αντικειμένων δέχονται ήδη τιμές διαφορετικών τύπων· η
        δυσκολία βρίσκεται στις συναρτήσεις τερματισμού, που πρέπει να γνωρίζουν τον τύπο
        κάθε στοιχείου για να το αφαιρέσουν. Μια λύση είναι η σταδιακή σύνθεση και
        αποσύνθεση της προσφοράς με μία κλήση ανά τύπο μέσα στο ίδιο \en{PTB}.
  \item \textbf{Πολλαπλοί στόχοι ανά προσφορά.} Αντικατάσταση του \code{target\_nft\_id}
        με λίστα αποδεκτών στόχων, ώστε ο ίδιος δεσμευμένος συνδυασμός να προσφέρεται για
        οποιοδήποτε από περισσότερα \en{NFTs}, με την πρώτη αποδοχή να καταναλώνει την
        προσφορά. Η απόρριψη θα αφαιρούσε μόνο τον στόχο του απορρίπτοντος από τη λίστα.
        Η ανάγκη προέκυψε από τη χρήση: με το τρέχον σχήμα, όποιος θέλει τα ίδια δύο
        \en{NFTs} για δύο διαφορετικούς στόχους πρέπει να διαλέξει έναν.
  \item \textbf{Περιοδική ανανέωση και φιλικότερα σφάλματα.} Ανανέωση των ερωτημάτων
        ανά τακτά διαστήματα, ώστε ενέργειες τρίτων να εμφανίζονται χωρίς αλλαγή σελίδας, και
        αντιστοίχιση των σφαλμάτων ανύπαρκτου αντικειμένου σε μήνυμα της μορφής «η προσφορά
        δεν είναι πλέον διαθέσιμη» (Ενότητα~\ref{sec:scenario-results}).
  \item \textbf{\en{Indexer} και διεπαφή \en{GraphQL}.} Αντικατάσταση της ανακάλυψης μέσω
        συμβάντων με ερωτήματα σε ευρετηριασμένη βάση, ώστε το κόστος ανάγνωσης να
        εξαρτάται από το πλήθος των ενεργών προσφορών. Τα συμβάντα που ήδη εκπέμπει το
        συμβόλαιο επαρκούν για την κατασκευή του ευρετηρίου.
  \item \textbf{Αυθαίρετοι τύποι στη διεπαφή.} Ανακάλυψη όλων των \en{NFTs} του χρήστη
        ανεξαρτήτως τύπου και απόδοσή τους μέσω του προτύπου \en{Display}, ώστε η εφαρμογή
        να χρησιμοποιείται με πραγματικές συλλογές του οικοσυστήματος.
  \item \textbf{Τυπική επαλήθευση.} Διατύπωση των αναλλοίωτων Α1 έως Α3 ως προδιαγραφών και
        απόδειξή τους με τον \en{Move Prover}, και ανεξάρτητος έλεγχος ασφάλειας, ως
        προϋποθέσεις για ανάπτυξη σε \en{mainnet}.
  \item \textbf{Ευκολότερη είσοδος χρηστών.} Υποστήριξη \en{zkLogin}, ώστε ένας νέος
        χρήστης να συνδέεται με λογαριασμό που ήδη διαθέτει αντί να εγκαθιστά πορτοφόλι,
        και χορηγούμενων συναλλαγών, ώστε να μην απαιτείται \en{SUI} για την πρώτη
        ενέργεια.
\end{enumerate}

Τέλος, αν και η εργασία αφορά \en{NFTs}, το πρωτόκολλο δεν κάνει καμία υπόθεση για το τι
αναπαριστούν τα αντικείμενα που ανταλλάσσονται. Με τον ίδιο κώδικα μπορούν να ανταλλαγούν
αντικείμενα παιχνιδιών, εισιτήρια, ή οποιοδήποτε περιουσιακό στοιχείο αναπαρίσταται ως
αντικείμενο του \en{Sui}.
